\documentclass[11pt]{article}
\usepackage[margin=1in]{geometry}
\usepackage{amsmath,amssymb}
\usepackage{natbib}
\usepackage{graphicx}
\usepackage{booktabs}
\usepackage{placeins}
\usepackage[font=small,labelfont=bf]{caption}
\usepackage[colorlinks=true,allcolors=blue]{hyperref}
\setlength{\emergencystretch}{3em}
\newcommand{\rr}{\langle r\rangle}

\title{\bfseries Narrow scalar clocks, not level repulsion:\\
a fourth spurious-GOE hazard and a re-audit of metallogenic timing statistics}
\author{Ruqing Chen\\[2pt]
\small GUT Geoservice Inc., Montr\'eal, Qu\'ebec, Canada\\
\small \texttt{ruqing.chen@gutgeoservice.com}}
\date{\today}

\begin{document}
\maketitle

\begin{abstract}
\noindent
A companion note documented three ways geophysical data fabricate Gaussian-ensemble
(GOE) ``level repulsion'' in the consecutive spacing ratio $\rr$ --- declustering,
seasonal gating, and the Wishart null. Here we add a fourth, and use it to
re-audit a published claim of temporal level repulsion in ore-forming systems. A
charge-and-release reservoir is a classical relaxation oscillator, i.e.\ a renewal
process; its inter-event intervals are independent, so it has no spectral rigidity
and cannot produce genuine level repulsion. Yet a renewal clock with a narrow
interval distribution reproduces \emph{any} value of $\rr$ --- we show it sweeps
the whole range from Poisson ($0.386$) through GOE ($0.531$) and GUE ($0.603$) up
to $0.88$ purely as a function of the interval coefficient of variation (CV). The
value of $\rr$ therefore carries no evidence of repulsion by itself. Re-analysing
the orogenic-gold geochronology of a published four-province study, we exactly
reproduce its reported $\rr=0.678$ and show that (i) it measures within-deposit
scatter among geochronometers, not province-scale timing; (ii) it is reproduced
by a zero-repulsion scalar clock of CV~$=0.38$; (iii) its apparent rigidity
(shuffle deficit $-0.031$) is an artifact of within-deposit normalisation, which a
zero-repulsion null reproduces identically; and (iv) it is unresolvable, since
jittering the ages within their stated $\pm5$~Myr uncertainties moves $\rr$ across
the entire Poisson-to-GUE range. We restate the program's organising principle
accordingly: a single isolated relaxation-oscillator source yields a scalar
renewal clock, not GOE level repulsion, while the superposition of many sources
yields Poisson; genuine Wigner--Dyson statistics require a chaotic, strongly-mixed
spectrum that a low-dimensional dissipative reservoir does not possess.
\end{abstract}

\section{Introduction}

The consecutive spacing ratio $\rr=\langle\min(s_i,s_{i+1})/\max(s_i,s_{i+1})\rangle$
\citep{oganesyan2007,atas2013} is an attractive, unfolding-free test for order in
a sequence, with reference values $0.386$ (Poisson), $0.531$ (GOE) and $0.603$
(GUE). A companion note \citep{chen2026spurious} showed that three mundane data
operations push $\rr$ to GOE-like values without any underlying repulsion:
minimum-separation declustering, seasonal gating of event times, and the
finite-sampling (Wishart) structure of correlation matrices. The unifying remedy
was the permutation (shuffle) test together with the coefficient of variation and
an appropriate null.

This note adds a fourth hazard that is both more elementary and, in geoscience
timing studies, more consequential: a \emph{narrow scalar clock}. Many proposed
``single-source'' Earth-system rhythms --- volcanic recharge, jökulhlaup
charge-and-release, ore-fluid focusing --- are physically relaxation oscillators:
a reservoir fills and discharges at a threshold, then refills. Such a process is a
\emph{renewal} process with independent inter-event intervals. It has, by
construction, no sequential correlation and no spectral rigidity, and so cannot
generate Wigner--Dyson level repulsion. Yet, as we show, a renewal clock with a
sufficiently narrow interval distribution reproduces any value of $\rr$, including
the GOE and GUE values, purely from its marginal. A reported $\rr\approx0.53$ is
therefore not evidence of level repulsion.

We make the point concrete by re-auditing a published study that interpreted
$\rr$-values of $0.57$--$0.71$ across four metallogenic provinces --- including a
non-magmatic orogenic-gold case --- as GOE/GUE ``level repulsion'' reflecting a
charge-and-release memory \citep{chen2026metallogeny}. The same program's spatial
companion \citep{chen2026sigma} had already cautioned that declustering fabricates
GOE; we show the temporal claim falls to the scalar-clock hazard and to the age
uncertainties, and we restate the organising principle so that the surviving,
defensible content is cleanly separated from the part that does not hold.

\section{The fourth hazard: a narrow scalar clock reproduces any \texorpdfstring{$\rr$}{r}}

\paragraph{Mechanism.} Model a charge-and-release source as a gamma renewal
process with mean-one intervals and coefficient of variation CV (smaller CV $=$
more regular firing). Figure~\ref{fig:mech}a shows $\rr$ as a function of CV: it
rises smoothly from the Poisson value at CV~$\approx1$ through GOE at
CV~$\approx0.52$ and GUE at CV~$\approx0.45$ to $0.88$ at CV~$=0.13$. A renewal
clock --- with no level repulsion whatsoever --- thus reproduces the entire range
of ``RMT'' $\rr$-values. The value of $\rr$ alone is uninformative about
repulsion; it is a one-to-one re-encoding of the interval CV.

\paragraph{Why this is not GOE.} Genuine Wigner--Dyson statistics arise from the
spectra of chaotic, strongly-mixed systems with the appropriate symmetry
\citep{bohigas1984}; integrable or near-integrable dynamics give Poisson
statistics \citep{berry1977}. A classical relaxation oscillator is the integrable
extreme --- a limit cycle --- and its event sequence is a renewal process, not a
Hamiltonian spectrum. It has no mechanism for level repulsion. The two genuine
fingerprints of GOE that a scalar clock lacks are (i) a small negative shuffle
deficit $\rr_{\rm obs}-\rr_{\rm shuf}<0$, from the sequential anti-correlation
(rigidity) of a repulsive spectrum, and (ii) a number variance that grows
logarithmically rather than linearly with interval length. Figure~\ref{fig:mech}b
shows the shuffle deficit: renewal clocks of every CV give
$\rr_{\rm obs}-\rr_{\rm shuf}\approx0$, while only the genuine GOE spectrum shows a
resolved negative value.

\paragraph{The resolution wall.} Crucially, that GOE signature is small
($\approx-0.03$) and requires a long sequence to detect: the standard error of
$\rr$ is $\approx\sigma_r/\sqrt{N}$ with $\sigma_r\approx0.27$, so resolving a
$-0.03$ deficit at $2\sigma$ needs $N\gtrsim300$ intervals. Geochronological
event sequences have tens. Below that, a scalar clock and a genuine GOE spectrum
are observationally indistinguishable in $\rr$, and any GOE claim is unsupported
regardless of the point value.

\begin{figure}[t]
\centering
\includegraphics[width=\textwidth]{fig_mechanism.pdf}
\caption{\textbf{(a)} A gamma renewal (charge-release) clock reproduces every
$\rr$ value --- Poisson, GOE, GUE and beyond --- purely as a function of its
interval coefficient of variation. \textbf{(b)} The discriminator: renewal clocks
of any CV give a shuffle deficit $\rr_{\rm obs}-\rr_{\rm shuf}\approx0$; only a
genuine GOE spectrum shows the resolved negative (rigidity) signature --- and even
that needs $N\gtrsim300$ levels to detect.}
\label{fig:mech}
\end{figure}

\section{Re-audit: orogenic-gold timing statistics}

We re-analyse the orogenic-gold dataset of \citet{chen2026metallogeny}, which
reported $\rr=0.678$ ($n=22$, ``best fit GUE'', Poisson excluded) and interpreted
it as intrinsic charge-release level repulsion in a non-magmatic ore system. The
analysis pipeline groups ages by deposit, forms the within-deposit spacings,
normalises them by each deposit's mean spacing, and pools the normalised spacings
across the ten deposits that have at least three dated ages.

We reproduce $\rr=0.678$ exactly, and find four problems
(Fig.~\ref{fig:real}).

\paragraph{(i) It measures the wrong quantity.} Each pooled ``interval'' is the
gap between two ages \emph{of the same deposit}, obtained from different
geochronometers (U--Pb xenotime, monazite, Re--Os, Ar--Ar). The statistic
therefore describes the scatter among dating methods on individual deposits, not
the province-scale spacing of distinct mineralising pulses that the charge-release
hypothesis is about. There is no province-scale temporal sequence in the
construction, so ``level repulsion'' --- a property of a sequence --- does not even
apply.

\paragraph{(ii) It is a scalar clock.} The pooled intervals have CV~$=0.31$, and a
zero-repulsion gamma renewal clock of CV~$=0.38$ reproduces $\rr=0.68$
(Section~2). The reported value is the marginal CV re-encoded; no repulsion is
present or required.

\paragraph{(iii) Its apparent rigidity is a normalisation artifact.} The shuffle
test gives $\rr_{\rm obs}-\rr_{\rm shuf}=-0.031$, superficially the GOE rigidity
signature. But dividing each deposit's intervals by their own mean turns them into
\emph{compositional} (closed) data: the normalised intervals of a deposit sum to a
fixed value (their count, $\sum_i s_i=\text{const}$), and this closure
algebraically forces a negative correlation between adjacent intervals. A pure
zero-repulsion renewal clock pushed through the \emph{same}
group-normalise-pool construction yields $\rr=0.69$ and a deficit of $-0.027$ ---
reproducing both the value and the ``rigidity'' (Fig.~\ref{fig:real}b). The
$-0.031$ is an artifact of the closure constraint, not spectral rigidity.

\paragraph{(iv) It is unresolvable.} The ages carry $\pm4$--$6$~Myr
uncertainties, comparable to the spacings. Jittering each age within its stated
error and re-running the exact pipeline moves $\rr$ across
$[0.39,0.67]$ --- the full Poisson-to-GUE range --- with the reported $0.678$
sitting at the upper edge of the noise (Fig.~\ref{fig:real}a). The published
bootstrap interval, which resamples intervals but ignores the age errors, is
correspondingly overconfident.

\begin{figure}[t]
\centering
\includegraphics[width=\textwidth]{fig_realdata.pdf}
\caption{\textbf{(a)} Re-running the exact pipeline with the ages jittered within
their $\pm$error moves $\rr$ across the whole Poisson-to-GUE range; the reported
$0.678$ is at the top edge. \textbf{(b)} A zero-repulsion scalar renewal clock
(CV~$=0.38$) put through the same group-normalise-pool construction reproduces
both the $\rr=0.68$ value and the $-0.03$ shuffle ``rigidity'': both are
construction artifacts, not level repulsion.}
\label{fig:real}
\end{figure}

\FloatBarrier
\section{Restated organising principle}

The four points are not a quirk of one dataset; they follow from the mechanism.
We therefore restate the program's organising principle
(Table~\ref{tab:principle}). A single isolated relaxation-oscillator source ---
the charge-and-release archetype --- produces a \emph{scalar renewal clock}: a
narrow interval distribution with high $\rr$ set by its CV, shuffle-invariant, and
without spectral rigidity. This is \emph{not} GOE level repulsion. Genuine
Wigner--Dyson statistics require a chaotic, strongly-mixed spectrum
\citep{bohigas1984}, which a fluid reservoir does not possess. The
superposition$\,\to\,$Poisson half of the principle is unaffected: pooling
independent sources broadens the interval distribution toward exponential and
drives $\rr\to0.386$, exactly as observed.

\begin{table}[t]
\centering
\caption{The organising principle, corrected. The high-$\rr$ ``single-source''
regime is a scalar clock, not level repulsion.}
\label{tab:principle}
{\small
\begin{tabular}{lll}
\toprule
Configuration & Signature & Interpretation\\
\midrule
Single relaxation oscillator   & high $\rr$, CV-set, shuffle-invariant & scalar renewal clock (\emph{not} GOE)\\
Chaotic / strongly-mixed spectrum & $\rr\!\approx\!0.53$, $\rr_{\rm obs}\!<\!\rr_{\rm shuf}$, rigid & GOE level repulsion\\
Superposition of many sources  & $\rr\to0.386$ & Poisson\\
\bottomrule
\end{tabular}}
\end{table}

This restatement strengthens rather than weakens the wider program, because it
places the ore-timing result in the same family as other correctly-diagnosed
scalar clocks --- jökulhlaup charge-release and steric nucleosome renewal
\citep{chen2026nucleosome} --- all of which show high, shuffle-invariant $\rr$ that
is marginal, not sequential.

\section{Implications for the audited studies}

\paragraph{Metallogenic timing \citep{chen2026metallogeny}.} The central claim ---
that single ore systems exhibit GOE/GUE level repulsion --- is not supported. The
high $\rr$-values are scalar clocks (and, for orogenic gold, unresolvable given
the age errors). What survives is the superposition result: pooling provinces
drives $\rr$ to the Poisson value, a genuine and correctly-interpreted effect. A
corrected version should drop the ``level repulsion'' interpretation and reframe
the single-source signal as a narrow scalar clock, retaining the
superposition$\,\to\,$Poisson contrast and the data compilation.

\paragraph{Sigma vein spacing \citep{chen2026sigma}.} The spatial results of that
study --- clustered, not stress-shadow-repulsive vein spacing, and the explicit
demonstration that declustering fabricates GOE --- are sound and unaffected. Only
its framing device, a ``time--space asymmetry'' contrasting temporal repulsion
with spatial clustering, rests on the temporal claim audited here. With the
temporal side reclassified as a scalar clock, the asymmetry should be restated as
``temporal scalar clock plus spatial clustering'', and the spatial and
methodological results stand on their own.

\section{Conclusion}

A narrow scalar (renewal) clock reproduces any value of the spacing ratio,
including the GOE and GUE values, with no level repulsion; in short geochronologic
sequences it is indistinguishable from a genuine repulsive spectrum, and age
uncertainties widen $\rr$ across the whole Poisson-to-GUE range. A published claim
of temporal GOE/GUE repulsion in ore systems is, on re-audit, a scalar clock
measured on within-deposit geochronometer scatter and unresolvable given the age
errors. The correct reading of a high, shuffle-invariant $\rr$ is a scalar clock,
not level repulsion --- the necessary fourth control to accompany the no-decluster,
de-seasonalise, and density-not-spacing rules of the companion note.

\section*{Data and code availability}
The orogenic-gold ages re-analysed here are from the dataset of
\citet{chen2026metallogeny}; a copy and all analysis code are at
\url{https://github.com/Ruqing1963/scalar-clock-not-goe} and archived on Zenodo
(DOI assigned on release). All other demonstrations are synthetic.

\begin{thebibliography}{99}

\bibitem[Ata\c{s} et al.(2013)]{atas2013}
Y.~Y. Ata\c{s}, E.~Bogomolny, O.~Giraud, and G.~Roux.
Distribution of the ratio of consecutive level spacings in random matrix
ensembles. \textit{Phys. Rev. Lett.} \textbf{110}, 084101 (2013).

\bibitem[Berry \& Tabor(1977)]{berry1977}
M.~V. Berry and M.~Tabor.
Level clustering in the regular spectrum.
\textit{Proc. R. Soc. Lond. A} \textbf{356}, 375--394 (1977).

\bibitem[Bohigas et al.(1984)]{bohigas1984}
O.~Bohigas, M.~J. Giannoni, and C.~Schmit.
Characterization of chaotic quantum spectra and universality of level fluctuation
laws. \textit{Phys. Rev. Lett.} \textbf{52}, 1--4 (1984).

\bibitem[Chen(2026a)]{chen2026sigma}
R.~Chen.
Time--space asymmetry in orogenic gold systems: temporal level repulsion but
spatial clustering of ore-stage veins at the Sigma deposit. Zenodo,
\url{https://doi.org/10.5281/zenodo.20778697} (2026).

\bibitem[Chen(2026b)]{chen2026metallogeny}
R.~Chen.
Charge-and-release rhythms in ore-forming systems: a random matrix theory test of
metallogenic timing across four provinces. Zenodo,
\url{https://doi.org/10.5281/zenodo.20768849} (2026).

\bibitem[Chen(2026c)]{chen2026nucleosome}
R.~Chen.
The consecutive spacing-ratio of nucleosome arrays: high regularity from steric
renewal, not level repulsion. Zenodo,
\url{https://doi.org/10.5281/zenodo.20839558} (2026).

\bibitem[Chen(2026d)]{chen2026spurious}
R.~Chen.
Spurious level repulsion in geophysical random-matrix analysis: declustering,
seasonal gating, and the Wishart null. Zenodo,
\url{https://doi.org/10.5281/zenodo.20883107} (2026).

\bibitem[Oganesyan \& Huse(2007)]{oganesyan2007}
V.~Oganesyan and D.~A. Huse.
Localization of interacting fermions at high temperature.
\textit{Phys. Rev. B} \textbf{75}, 155111 (2007).

\end{thebibliography}

\end{document}
